\section{xLSTM Model Architecture}

The xLSTM implementation provided by NXAI uses the \verb|xlstm| package's \verb|xLSTMLMModel| as an autoregressive language model over REMI+ token IDs. The training code fills in \verb|vocab_size| from the tokenizer and \verb|context_length| from the data block size at runtime. Therefore, the model YAML files specify architectural scaling parameters rather than corpus-specific values.

The repository defines a simple scaling ladder in \verb|configs/model/xlstm/|:
\begin{itemize}
	\item \verb|basic|: embedding dimension 128, 4 blocks, one sLSTM block;
	\item \verb|small|: embedding dimension 192, 6 blocks, two sLSTM blocks;
	\item \verb|base|: embedding dimension 256, 8 blocks, three sLSTM blocks;
	\item \verb|medium|: embedding dimension 384, 10 blocks, three sLSTM blocks and 8 heads.
\end{itemize}

The common design keeps the general recipe fixed and scales mainly through embedding dimension and number of blocks. The mLSTM and sLSTM sub-blocks use a 1D convolution kernel size of 4. The small/base configurations use 4 heads; the medium configuration uses 8 heads. sLSTM feedforward layers use a projection factor of 1.3 and GELU activation. The CUDA sLSTM backend is used in the GPU-oriented configs, while a CPU-compatible basic config exists separately.

\begin{table}[h]
	\centering
	\caption{xLSTM scaling ladder from \texttt{configs/model/xlstm/README.md}.}
	\label{tab:xlstm_scale}
	\begin{tabular}{lrrrr}
		\toprule
		Config & Embedding dim. & Blocks & sLSTM positions & Approx. params \\
		\midrule
		basic  & 128            & 4      & [1]             & 0.64M          \\
		small  & 192            & 6      & [1, 4]          & 1.70M          \\
		base   & 256            & 8      & [1, 3, 5]       & 3.77M          \\
		medium & 384            & 10     & [1, 4, 7]       & 9.39M          \\
		\bottomrule
	\end{tabular}
\end{table}

The current completed experiment summaries show 1,628,384 trainable parameters for the \verb|small| run and 3,671,080 trainable parameters for the \verb|base| runs with tokenizer 11. These values differ slightly from the rounded README estimates because the actual vocabulary size and runtime configuration affect the final language-model head.

